% ═══════════════════════════════════════════════════════════════
% ARBEITSBLATT: Kernreaktionen (Physik Klasse 10)
% Begleitend zum digitalen Lernpfad
% ═══════════════════════════════════════════════════════════════

\documentclass[10pt, a4paper]{article}

\usepackage[utf8]{inputenc}
\usepackage[T1]{fontenc}
\usepackage[ngerman]{babel}
\usepackage{helvet}
\renewcommand{\familydefault}{\sfdefault}

\usepackage[margin=1.5cm, top=1.8cm, bottom=1.5cm]{geometry}
\usepackage{enumitem}
\usepackage{tabularx}
\usepackage{booktabs}
\usepackage{array}
\usepackage{multicol}

\usepackage{xcolor}
\usepackage{tcolorbox}
\usepackage{amssymb}
\usepackage{fancyhdr}
\usepackage{lastpage}

\setlist{nosep, leftmargin=1.5em}

% Farben
\definecolor{physik}{HTML}{2c5aa0}
\definecolor{grau}{HTML}{888888}
\definecolor{hellgrau}{HTML}{f0f0f0}
\definecolor{fehler}{HTML}{c62828}
\definecolor{fehlerbg}{HTML}{ffebee}

\newcommand{\fachfarbe}{physik}

% Variablen
\newcommand{\thema}{Kernreaktionen}
\newcommand{\fach}{Physik}
\newcommand{\klasse}{Klasse 10}
\newcommand{\stunde}{Stunde 7-8}

% Kopf- und Fußzeile
\pagestyle{fancy}
\fancyhf{}
\renewcommand{\headrulewidth}{0.4pt}
\lhead{\small\fach{} \klasse{} -- \thema}
\rhead{\small Name: \rule{4cm}{0.4pt}}
\cfoot{\small\thepage/\pageref{LastPage}}

% Befehle
\newcommand{\luecke}[1]{\rule{#1}{0.4pt}}

\newtcolorbox{merksatz}[1][]{
    colback=hellgrau, colframe=\fachfarbe, boxrule=0.8pt, arc=0pt,
    left=6pt, right=6pt, top=6pt, bottom=6pt,
    fonttitle=\bfseries\small, title=#1
}

\newtcolorbox{fehlerbox}{
    colback=fehlerbg, colframe=fehler, boxrule=0.8pt, arc=0pt,
    left=6pt, right=6pt, top=4pt, bottom=4pt,
    fonttitle=\bfseries\small, title=Typische Fehler
}

\newcommand{\abschnitt}[1]{%
    \vspace{8pt}\noindent\textbf{\textcolor{\fachfarbe}{#1}}\vspace{4pt}
}

\begin{document}

\begin{center}
{\large\bfseries Arbeitsblatt: \thema{} -- \stunde}\\[2pt]
{\small\textcolor{grau}{Begleitend zum Lernpfad -- in den Hefter übernehmen}}
\end{center}
\vspace{-2pt}\hrule\vspace{8pt}

\begin{multicols}{2}
\small

% ─────────────────────────────────────────────────────────────
% KASTEN 1: Kernspaltung
% ─────────────────────────────────────────────────────────────

\abschnitt{Kasten 1: Kernspaltung}

\begin{merksatz}[Definition und Beispiel]
Bei der \textbf{Kernspaltung} wird ein schwerer Kern

durch ein \luecke{2cm} in zwei mittelschwere

Kerne gespalten.

\vspace{6pt}

\textbf{Beispiel U-235:}

$^{235}_{92}$U + $^{1}_{0}$n $\rightarrow$ \luecke{1.5cm} + \luecke{1.5cm} + 3n

\vspace{6pt}

Dabei wird \luecke{2cm} freigesetzt.
\end{merksatz}

\vspace{6pt}

% ─────────────────────────────────────────────────────────────
% KASTEN 2: Kettenreaktion
% ─────────────────────────────────────────────────────────────

\abschnitt{Kasten 2: Kettenreaktion}

\begin{merksatz}[Kontrollierte vs. unkontrollierte]
\textbf{Kettenreaktion:}

Die bei der Spaltung freiwerdenden \luecke{2cm}

können weitere Kerne spalten.

\vspace{6pt}

\renewcommand{\arraystretch}{1.4}
\begin{tabular}{|p{2.5cm}|p{3.5cm}|}
\hline
\textbf{Art} & \textbf{Anwendung} \\
\hline
Kontrolliert & \luecke{2.5cm} \\
\hline
Unkontrolliert & \luecke{2.5cm} \\
\hline
\end{tabular}
\renewcommand{\arraystretch}{1}

\vspace{6pt}

\textbf{Kritische Masse:} Mindestmenge, ab der

eine Kettenreaktion \luecke{2.5cm}.
\end{merksatz}

\vspace{6pt}

% ─────────────────────────────────────────────────────────────
% KASTEN 3: Kernfusion
% ─────────────────────────────────────────────────────────────

\abschnitt{Kasten 3: Kernfusion}

\begin{merksatz}[Definition und Beispiel]
Bei der \textbf{Kernfusion} verschmelzen

\luecke{2cm} Kerne zu einem schwereren Kern.

\vspace{6pt}

\textbf{Beispiel (Sonne):}

$^{2}_{1}$H + $^{3}_{1}$H $\rightarrow$ \luecke{1.5cm} + $^{1}_{0}$n + Energie

\vspace{6pt}

Voraussetzung: Extrem hohe \luecke{2cm}

(ca. 100 Millionen °C)
\end{merksatz}

\columnbreak

% ─────────────────────────────────────────────────────────────
% ÜBUNG 1: Zuordnung (AFB I)
% ─────────────────────────────────────────────────────────────

\abschnitt{Übung 1: Spaltung oder Fusion?}

Ordne zu: Spaltung (S) oder Fusion (F)

\vspace{4pt}
\begin{tabular}{|p{4.5cm}|c|}
\hline
\textbf{Aussage} & \textbf{S/F} \\
\hline
Energiequelle der Sonne & \\
\hline
Nutzt Uran-235 & \\
\hline
Verschmelzung leichter Kerne & \\
\hline
Funktioniert im Kernkraftwerk & \\
\hline
Benötigt 100 Mio. °C & \\
\hline
Produziert radioaktiven Müll & \\
\hline
\end{tabular}

\vspace{8pt}

% ─────────────────────────────────────────────────────────────
% ÜBUNG 2: Gleichungen (AFB II)
% ─────────────────────────────────────────────────────────────

\abschnitt{Übung 2: Reaktionsgleichungen}

\textbf{a)} Vervollständige die Kernspaltung:

$^{235}_{92}$U + $^{1}_{0}$n $\rightarrow$ $^{92}_{36}$Kr + \luecke{1.5cm} + 3$^{1}_{0}$n

\vspace{6pt}

\textbf{b)} Berechne: Wie viele Neutronen entstehen

nach 4 Generationen der Kettenreaktion,

wenn pro Spaltung 3 Neutronen frei werden?

\vspace{4pt}
\luecke{5.8cm}

\vspace{8pt}

% ─────────────────────────────────────────────────────────────
% ÜBUNG 3: Fehler erklären (AFB III)
% ─────────────────────────────────────────────────────────────

\abschnitt{Übung 3: Fehler erklären (AFB III)}

\begin{fehlerbox}
Erkläre, warum diese Aussage falsch ist:

\vspace{4pt}

\textit{„Ein Kernkraftwerk kann wie eine Atombombe explodieren."}

\vspace{4pt}
\luecke{5.8cm}

\vspace{4pt}
\luecke{5.8cm}
\end{fehlerbox}

\vspace{8pt}

% ─────────────────────────────────────────────────────────────
% ÜBUNG 4: Transfer (AFB III)
% ─────────────────────────────────────────────────────────────

\abschnitt{Übung 4: Vergleich}

Nenne je zwei Vor- und Nachteile von Kernspaltung und Kernfusion:

\vspace{4pt}
\luecke{5.8cm}

\vspace{4pt}
\luecke{5.8cm}

\vspace{4pt}
\luecke{5.8cm}

\end{multicols}

% ═══════════════════════════════════════════════════════════════
% SEITE 2: SIMULATION KETTENREAKTION
% ═══════════════════════════════════════════════════════════════

\newpage

\begin{center}
{\large\bfseries Simulationsprotokoll: Kettenreaktion}\\[2pt]
{\small\textcolor{grau}{PhET-Simulation: Nuclear Fission}}
\end{center}
\vspace{-2pt}\hrule\vspace{12pt}

\abschnitt{Teil A: Einzelne Spaltung beobachten}

\begin{merksatz}[Beobachtung]
Beschreibe, was bei einer einzelnen Kernspaltung passiert:

\vspace{4pt}
1. Was trifft auf den U-235-Kern? \luecke{6cm}

\vspace{4pt}
2. Was passiert mit dem Kern? \luecke{6cm}

\vspace{4pt}
3. Was wird freigesetzt? \luecke{6cm}
\end{merksatz}

\vspace{8pt}

\abschnitt{Teil B: Kettenreaktion simulieren}

Starte mit einem Neutron und zähle die Spaltungen pro Generation:

\vspace{6pt}

\begin{tabular}{|c|c|c|c|}
\hline
\textbf{Generation} & \textbf{Neutronen} & \textbf{Spaltungen} & \textbf{Beobachtung} \\
\hline
1 & 1 & \luecke{1cm} & \luecke{4cm} \\
\hline
2 & \luecke{1cm} & \luecke{1cm} & \luecke{4cm} \\
\hline
3 & \luecke{1cm} & \luecke{1cm} & \luecke{4cm} \\
\hline
4 & \luecke{1cm} & \luecke{1cm} & \luecke{4cm} \\
\hline
\end{tabular}

\vspace{12pt}

\abschnitt{Teil C: Kontrollstäbe testen}

\textbf{1. Starte eine Kettenreaktion OHNE Kontrollstäbe.}

Was passiert? \luecke{10cm}

\vspace{8pt}

\textbf{2. Starte erneut MIT Kontrollstäben.}

Was passiert? \luecke{10cm}

\vspace{8pt}

\textbf{3. Erkläre die Funktion der Kontrollstäbe:}

\luecke{14cm}

\vspace{4pt}
\luecke{14cm}

\vspace{12pt}

\abschnitt{Teil D: Auswertung}

\textbf{1. Warum ist die kritische Masse wichtig?}

\vspace{4pt}
\luecke{14cm}

\vspace{8pt}

\textbf{2. Was unterscheidet einen Kernreaktor von einer Atombombe? (Nennt 2 Unterschiede)}

\vspace{4pt}
\luecke{14cm}

\vspace{4pt}
\luecke{14cm}

\vspace{8pt}

\textbf{3. Warum ist Kernfusion auf der Erde so schwer zu realisieren?}

\vspace{4pt}
\luecke{14cm}

\vspace{4pt}
\luecke{14cm}

\end{document}
